% config/hyphenation.tex
% -----------------------------------------------------------------------------
% Hyphenation (pemenggalan kata) untuk bahasa Indonesia.
%
% LaTeX memiliki pola hyphenation untuk bahasa Indonesia (via babel),
% tetapi kamusnya belum selengkap bahasa Inggris.
% Jika ada kata yang tidak terpotong dengan benar (overshoot margin),
% tambahkan kata tersebut di sini dengan tanda hubung (-) di tiap suku kata.
%
% Contoh:
%   \hyphenation{me-nya-ya-ngi pe-ne-li-ti-an pe-ngem-ba-ngan}
%
% Petunjuk pemenggalan bahasa Indonesia (EYD):
%   1. Pemenggalan dilakukan pada batas suku kata.
%   2. Dua vokal berurutan (diftong) tidak dipisah: au, ai, ei, oi.
%   3. Konsonan di antara dua vokal ikut ke suku kata berikutnya: ba-ha-sa.
%   4. Dua konsonan berurutan dipisah di tengah: kon-so-nan.
%   5. Awalan, akhiran, dan imbuhan dipisah jika membentuk suku kata.
%
% Lihat https://ejaan.kemdikbud.go.id/eyd/ untuk panduan lengkap.
% -----------------------------------------------------------------------------

\hyphenation{
  % Tambahkan kata-kata di sini, pisahkan dengan spasi
  ana-li-sa
  apa-li-ka-si
  dae-rah
  di-guna-kan
  eks-pe-ri-men
  iden-ti-fi-ka-si
  imple-men-ta-si
  in-for-ma-si
  kom-pu-ter
  kon-sep
  kon-so-nan
  kor-pus
  me-to-de
  pe-ngem-ba-ngan
  pe-ne-li-ti-an
  pe-nga-lam-an
  per-ban-ding-an
  pe-ran-cang-an
  per-ma-sa-la-han
  pre-fe-rens
  pro-gram
  pro-gram stu-di
  re-fe-ren-si
  re-ko-men-da-si
  tek-no-lo-gi
  ter-se-but
  va-li-da-si
}
